\section{Algebraische Gesetze}
\label{sec:fundamentals_algebraic_laws}

Für die bekannten Zahlenmengen gelten drei grundlegende Gesetze der Algebra, die Ihnen vertraut sein sollten. \\

\noindent\begin{minipage}{\linewidth}
\begin{DefinitionBox}{Kommutativgesetz}{fundamentals_commutative_law}
  Das \emph{Kommutativgesetz} \EnglishTerm{commutative law} besagt, dass die
  Reihenfolge der Operanden $a$ und $b$ bei der Addition und der Multiplikation
  keine Rolle spielt:
  \begin{align*}
    a+b &= b+a, \\
    a\cdot b &= b\cdot a.
  \end{align*}
  Daher wird das Kommutativgesetz auch Vertauschungsgesetz genannt.
\end{DefinitionBox}

\begin{ExampleBox}{Kommutativgesetz}{fundamentals_commutative_law_example}
  \begin{enumerate}[label=(\alph*),itemsep=0.6em]
    \item $3+2=\Cloze{3.5cm}{5=2+3}$
    \item $3\cdot 2=\Cloze{3.5cm}{6=2\cdot 3}$
  \end{enumerate}
\end{ExampleBox}
\end{minipage}

\vspace{1cm}

\noindent\begin{minipage}{\linewidth}
\begin{DefinitionBox}{Assoziativgesetz}{fundamentals_associative_law}
  Das \emph{Assoziativgesetz} \EnglishTerm{associative law} besagt, dass die
  Gruppierung der Operanden $a$, $b$ und $c$ durch Klammern keinen Einfluss auf
  das Ergebnis hat:
  \begin{align*}
    (a+b)+c &= a+(b+c), \\
    (a\cdot b)\cdot c &= a\cdot(b\cdot c).
  \end{align*}
  Werden viele Operanden addiert oder multipliziert, darf frei gewählt werden,
  welches Operandenpaar zuerst verrechnet wird.
\end{DefinitionBox}

\begin{ExampleBox}{Assoziativgesetz}{fundamentals_associative_law_example}
  \begin{enumerate}[label=(\alph*),itemsep=0.8em]
    \item \emph{Addition:}
      \[
        (2+3)+4=\Cloze{3.8cm}{5+4=9}
        \qquad\text{und}\qquad
        2+(3+4)=\Cloze{3.8cm}{2+7=9}.
      \]
    \item \emph{Multiplikation:}
      \[
        (2\cdot 3)\cdot 4=\Cloze{3.8cm}{6\cdot 4=24}
        \qquad\text{und}\qquad
        2\cdot(3\cdot 4)=\Cloze{3.8cm}{2\cdot 12=24}.
      \]
  \end{enumerate}
  Beide möglichen Gruppierungen führen jeweils zum selben Ergebnis!
\end{ExampleBox}
\end{minipage}

\vspace{1cm}

\noindent\begin{minipage}{\linewidth}
Während das Kommutativ- und das Assoziativgesetz jeweils für die Addition und
die Multiplikation gelten, verbindet das Distributivgesetz beide Rechenoperationen.
%
\begin{DefinitionBox}{Distributivgesetz}{fundamentals_distributive_law}
  Das \emph{Distributivgesetz} \EnglishTerm{distributive law} verbindet die
  Multiplikation mit der Addition:
  \[
    a\cdot(b+c)=a\cdot b+a\cdot c.
  \]
  Es erlaubt das Ausmultiplizieren von Klammern und umgekehrt das Ausklammern
  gemeinsamer Faktoren aus Summen. Beim Auflösen der Klammer wird jeder
  Summand in der Klammer mit dem Faktor vor der Klammer multipliziert.
\end{DefinitionBox}

\begin{ExampleBox}{Distributivgesetz}{fundamentals_distributive_law_example}
  \begin{enumerate}[label=(\alph*),itemsep=0.8em]
    \item \emph{Ausmultiplizieren:}
      \[
        a\cdot\left(5+b+c^2\right)
        =\Cloze{7cm}{5a+ab+ac^2}.
      \]
    \item \emph{Ausklammern:}
      \[
        a\cdot b\cdot c^2+a\cdot b
        =\Cloze{9cm}{a\cdot b\cdot c^2+a\cdot b\cdot 1
          =a\cdot b\left(c^2+1\right)}.
      \]
      Wird beim Ausklammern aus einem Summanden das vollständige Produkt herausgezogen, bleibt in der Klammer eine $1$ zurück.
  \end{enumerate}
\end{ExampleBox}
\end{minipage}

\vspace{1cm}
Wenn wir später mit anderen Objekten wie komplexen Zahlen, Vektoren oder Matrizen rechen,
werden manche dieser 3 Gesetzte nicht mehr gültig sein.
